\section{The high-prime local tail and the exact remaining gates}
\label{sec:gates}

We isolate only the local Euler tail contributed by individual primes
\(p>R\).  This is not the whole arithmetic tail of the sharp
\(q\le R\) Ramanujan cutoff: composite squarefree moduli \(q>R\)
built entirely from primes at most \(R\) are also omitted and are not
controlled by the theorem below.  Let \(m_1\ne m_2\), put
\begin{equation}
 A=m_1m_2(m_1-m_2),
 \label{eq:A-discriminant}
\end{equation}
and, for each prime \(p\), let
\begin{equation}
 \nu_p=\#\{j\bmod p:(m_1j+h)(m_2j+h)\equiv0\pmod p\}.
 \label{eq:nu-p}
\end{equation}

\begin{theorem}[Elementary high-prime Euler tail]
\label{thm:euler-tail}
Assume \(R>\max(2,|h|)\).  Then
\begin{equation}
 \left|
 \log\prod_{p>R}
 \frac{1-\nu_p/p}{(1-1/p)^2}
 \right|
 \ll
 \frac1R+
 \frac{\log(2|A|)}{R\log R}.
 \label{eq:euler-tail}
\end{equation}
The implied constant is absolute.
\end{theorem}

\begin{proof}
Since \(p>R>|h|\), no such prime divides \(h\).  If \(p\nmid A\),
the two affine forms have two distinct roots, so \(\nu_p=2\), and
\begin{equation}
 \log\frac{1-2/p}{(1-1/p)^2}=O(p^{-2}).
 \label{eq:generic-euler-factor}
\end{equation}
If \(p\mid A\), then \(0\le\nu_p\le2\) and the same logarithm is
\(O(p^{-1})\).  Hence the left side of \eqref{eq:euler-tail} is at
most
\[
 O\!\left(\sum_{p>R}p^{-2}\right)
 +O\!\left(\sum_{\substack{p>R\\p\mid A}}p^{-1}\right).
\]
The first sum is \(O(R^{-1})\).  For the second, use
\[
 \frac1p\le\frac{\log p}{R\log R}
 \quad(p>R),
 \qquad
 \sum_{p\mid A}\log p\le\log|A|,
\]
with \(\log(2|A|)\) covering \(|A|=1\).
\end{proof}

\begin{remark}[A local statement only]
\Cref{thm:euler-tail} says that individual primes above \(R\) do not
hide a large unrecorded local Euler factor.  It says nothing about the
omitted composite moduli just described, and it is not an estimate for a
finite \(j\)-correlation, a long-lcm Poisson remainder, a prime-pair
sum, or any of the four locally calibrated channels.  In particular,
it does not
replace determinant dispersion.
\end{remark}

We can now state the remaining target without ambiguity.  Let
\(\mathscr G\) be the generic row set in \eqref{eq:generic-rows} and
let \(\cC_D(\mathscr G)\) be \eqref{eq:abstract-correlation} with the
actual TPC-18 coefficients.

\begin{definition}[Primitive residual-dispersion target]
\label{def:residual-dd}
For a prescribed \(B>0\), the target estimate is
\begin{equation}
 |\cC_D(\mathscr G)|
 \ll_{B,h,\delta,\kappa}XQ(\log X)^{-B}
 \label{eq:residual-dd-target}
\end{equation}
uniformly in the surviving dyadic parameters.  This definition is a
named missing estimate, not a theorem of the present paper.
\end{definition}

The exact identities above split attempts at
\eqref{eq:residual-dd-target} into two auditable routes.

\begin{description}[leftmargin=2.1em,labelindent=0em,style=nextline]
\item[Direct determinant-frequency route.]
Derive a common Kuznetsov/dispersion form for the four terms in
\eqref{eq:centered-four-channel}, before absolute values, so that their
local main spectra cancel.  For the resulting pair-sequence measure,
verify the precise hypotheses and concentration functional of
\cite[Theorem~5.2]{Pascadi2026LargeSieve}, as described in
\cref{rem:mobius-concentration}.  The explicit matched reduction and
the sharp concentration estimate are both missing.

\item[CRT--Poisson long-lcm route.]
Apply the all-modulus residual CRT expansion termwise, or first derive
an explicitly matched four-channel analogue.  Evaluate every
short-period zero mode and bound the complete complementary long
spectrum while keeping the joint, two single, and constant terms
together.  \Cref{cor:whole-model-completion} supplies only the finite
\(MM\) subcalculation.  The missing long estimate must tolerate the
coupled \(g,a,b\) moduli, the determinant constraint, and the sharp
source--M\"obius coefficients typified by
\eqref{eq:three-inverse-phases}.
\end{description}

These two matched-spectrum routes are sufficient, not logically
exhaustive.  A third possibility would be a separate global estimate
for each channel that exploits cancellation over the full row family.
\Cref{thm:model-witness} says only that such an estimate cannot come
from rowwise complete-period centering of the \(MM\) term.

\begin{remark}[Route ledger, not a closure theorem]
The exact reductions identify the following sufficient future
packages:
\begin{enumerate}[label=\textnormal{(\roman*)}]
 \item an explicitly derived matched four-channel spectral reduction,
 followed by verification of the precise concentration inequality
 required for its resulting measure, with total error
 \(O_B(XQ(\log X)^{-B})\);
 \item an all-modulus residual or matched four-channel Poisson
 decomposition, exact evaluation of all short zero modes, and a bound
 \(O_B(XQ(\log X)^{-B})\) for the complete long spectrum.
\end{enumerate}
If a future work supplies either package with the displayed total
error, then \cref{thm:drift-small,thm:residual-crt} show that the
primitive block reaches \eqref{eq:residual-dd-target}.  This is a
sufficiency ledger, not a quantitatively stated conditional result:
neither package is supplied here.  \Cref{thm:euler-tail} separately
records only that the local-factor ledger has no large omitted
high-prime tail; it supplies no correlation estimate.
\end{remark}

\begin{remark}[What closing this ledger would mean]
Even a proof of \eqref{eq:residual-dd-target} closes only the primitive
residual block delivered by TPC-18.  A twin-prime theorem would still
require every preceding decomposition, endpoint, positivity, and
parameter-uniformity step in the full chain.  Conversely, a proof that
a precisely stated concentration inequality fails in the required
parameter range, or construction of a long-lcm packet that
quantitatively violates the target bound, would be a rigorous
obstruction result, not evidence for or against the twin-prime
conjecture.
\end{remark}
